% Generated draft: requires mathematical and source review.
\documentclass[11pt]{article}
\usepackage[margin=25mm]{geometry}
\usepackage{fontspec}
\setmainfont{Latin Modern Roman}
\usepackage{amsmath,amssymb,mathtools,unicode-math}
\setmathfont{Latin Modern Math}
\usepackage{enumitem,needspace,xurl,hyperref}
\setlength{\parindent}{0pt}
\setlength{\parskip}{5pt}
\setlength{\emergencystretch}{3em}
\newcommand{\N}{\mathbb{N}}
\newcommand{\Z}{\mathbb{Z}}
\newcommand{\Q}{\mathbb{Q}}
\newcommand{\R}{\mathbb{R}}
\newcommand{\C}{\mathbb{C}}
\newcommand{\F}{\mathbb{F}}
\newcommand{\A}{\mathbb{A}}
\newcommand{\PP}{\mathbb{P}}
\newcommand{\E}{\mathbb{E}}
\newcommand{\eps}{\varepsilon}
\newcommand{\dd}{\,\mathrm{d}}
\begin{document}
\section*{UnsolvedMath ID: 20002480}
\textbf{Source problem code:} catalogue entry
\subsection{Definitions and mathematical statement}
Put $\mathbb T^n=\R^n/\Z^n$. The question attributed to Kostia Khanin concerns a smooth Tonelli Lagrangian $L_0(x,v)$ on $\mathbb T^n\times\R^n$, convex and superlinear in $v$ with positive-definite velocity Hessian, formally perturbed by
\[
 L(x,v,t)=L_0(x,v)+\sum_{i=1}^{N}F_i(x)\dot w_i(t),
\]
where the $w_i$ are independent Brownian motions. The natural forcing vector is $F=(F_1,\ldots,F_N):\mathbb T^n\to\R^N$, assumed to be an embedding. The extract instead writes $F_1,\ldots,F_n$ and an undefined target $M$ for the minimizer; the aid explicitly repairs these to $N$ components and $\gamma:\R\to\mathbb T^n$. Since $\dot w_i$ is distributional, the action and random flow require a pathwise stochastic interpretation, not an ordinary smooth time-dependent Lagrangian.

A global minimizer means a curve whose restriction to every finite interval minimizes the pathwise action among curves with the same endpoints. The source states almost-sure uniqueness and conjectures almost-sure hyperbolicity of this minimizer; it attributes the one-dimensional result to E--Khanin--Mazel--Sinai. It does not fully specify cohomology/tilt, action convention, flow regularity, or whether hyperbolicity means a nonuniform Lyapunov splitting or a uniform exponential dichotomy.

The aid's special case assumes a separable system
\[
 L_0(x,v)=\sum_{j=1}^n\ell_j(x_j,v_j),\qquad
 \sum_iF_i(x)\dot w_i(t)=
 \sum_{j=1}^n\sum_{k=1}^{N_j}f_{j,k}(x_j)\dot w_{j,k}(t),
 \quad N=\sum_jN_j.
\]
Assume every one-dimensional factor has, almost surely, a unique global minimizer $\gamma_j$ and a derivative cocycle with Lyapunov exponents $-\lambda_j,\lambda_j$, $\lambda_j>0$, under the requisite integrability hypotheses. Then the reported product result is that the global minimizer is uniquely $\gamma=(\gamma_1,\ldots,\gamma_n)$ and
\[
 D\Phi_t=\bigoplus_{j=1}^nD\Phi_{j,t},\qquad
 \operatorname{Lyap}(D\Phi)=\{-\lambda_1,\lambda_1,\ldots,-\lambda_n,\lambda_n\}
\]
with multiplicities. In particular its gap from zero is exactly $\min_j\lambda_j>0$. The assertions hold on the intersection of the finitely many factor probability-one events. If each map $(f_{j,1},\ldots,f_{j,N_j}):\mathbb T\to\R^{N_j}$ is an embedding, their coordinate-block product is an embedding of $\mathbb T^n$, preserving the stated forcing hypothesis.

The action splitting and block cocycle justify this separable higher-dimensional case. They do not address general coupled $L_0$ or forcing. The spectral-gap statement concerns Lyapunov hyperbolicity; a stronger uniform-in-time dichotomy would require corresponding factor assumptions and estimates.
\subsection{Short English statement}
The conjecture asks whether the unique random action minimizer is almost surely hyperbolic in higher dimensions. The supplied result handles independent coordinate blocks: minimizers multiply, Lyapunov spectra combine, and the smallest positive factor exponent is the product gap. Coupled systems remain outside this argument.
\subsection{Sources}
\begin{itemize}
\item Local UnsolvedMath registry, permanent ID 20002480, source code catalogue entry; created and updated 14 August 2026. Canonical \texttt{aim-pdes-notes.json}, zero-based \texttt{notes[91]}, tag \texttt{problem}; extracted/background item 2, attributed to Kostia Khanin.
\item AIM ``New connections between dynamical systems and PDE's,'' version dated 15 August 2003: \url{https://aimath.org/WWN/dynpde/dynpde.pdf}. The supplied verification identifies the official problem as 22 rather than the extracted 2. The registry marks the statement reconstructed\_unverified. Original remarks and literature field are empty.
\item Research attempt 1: \texttt{partial\_result}, problem status at run \texttt{partially\_solved}, rigor label \texttt{proved}, novelty confidence low. The one-dimensional theorem is attributed to E--Khanin--Mazel--Sinai without a full title, date, theorem number, or authoritative bibliographic link.
\item Ulam AI / UnsolvedMath Contributors, CC BY 4.0. Machine-generated research aid, not peer reviewed; independently verify before citation.
\end{itemize}
\end{document}
